% =============================================================================
% IEEE Standard Format — Assignment 2
% =============================================================================
\documentclass[conference]{IEEEtran}

% --- Packages ---
\usepackage[utf8]{inputenc}
\usepackage[T1]{fontenc}
\usepackage{graphicx}
\usepackage{booktabs}
\usepackage{amsmath}
\usepackage{amssymb}
\usepackage{hyperref}
\usepackage{xcolor}
\usepackage{float}
\usepackage{enumitem}
\usepackage{cite}
\usepackage{url}
\usepackage{balance}

\graphicspath{{result/}}

\hypersetup{colorlinks=true,linkcolor=blue!70!black,citecolor=green!50!black,urlcolor=blue!60!black}

% =============================================================================
\begin{document}

% --- Title ---
\title{Evaluating the Bias, Trustworthiness and Fairness of Open-Source Large Language Models: A DECODINGTRUST Analysis of Phishing Susceptibility}

\author{
\IEEEauthorblockN{[Student Name]}
\IEEEauthorblockA{Student ID: [Student ID]\\
Advanced Topics in Artificial Intelligence and Machine Learning\\
Assignment 2 --- Individual (30\%)\\
Email: [student@email.com]}
}

\maketitle

% --- Abstract ---
\begin{abstract}
The use of large language models (LLMs) in security settings has grown steadily, but few studies have examined the effect of their demographic bias tendencies when gauging human vulnerability. This study builds upon the DECODINGTRUST~\cite{wang2023decodingtrust} benchmark and analyzes the performance of five publicly available large language models: Qwen~2.5, DeepSeek~R1, Microsoft Phi-2, Mistral~7B, and TinyLlama using seven trust dimensions. More than 1,000 phishing examples involving personas were produced using the two prompts evaluation setup method for the following seven dimensions: toxicity, stereotype bias, adversarial robustness, privacy, factuality, fairness, and machine ethics. The findings show that some of these models drastically change their vulnerability evaluations according to only the age or neurodivergence of a persona. Model performance is visualized using a radar graph to analyze the effect of model size and instruction tuning on model safety and fairness.
\end{abstract}

\begin{IEEEkeywords}
Large Language Models, Phishing Susceptibility, Bias, Fairness, DECODINGTRUST, Security
\end{IEEEkeywords}

% =============================================================================
% SECTION 1
% =============================================================================
\section{Introduction}

Phishing continues to be among the leading cyber threats, with 4.7~million attacks recorded during 2022 alone~\cite{apwg2022}. Recent advances in generative AI have further complicated this landscape, enabling automated social engineering at scale. As firms start to adopt LLMs to automate security evaluations and assess phishing vulnerability~\cite{bommasani2021opportunities}, a key ethical dilemma emerges: are these systems unbiased, or do they depend on stereotyping based on demographic factors to determine vulnerability?

A firm cybersecurity platform powered by an LLM may discriminate against certain groups of workers, e.g., older employees and people with ADHD, thus discriminating against them. Such an action will lead to the creation of workplace disparity while collecting inaccurate security data.

This project focuses on answering three fundamental questions:
\begin{enumerate}[leftmargin=*]
 \item Is there any use of demographic stereotypes (e.g., age, gender, neurodivergence) in deciding the vulnerability of the target user in the evaluated open-source LLMs?
 \item Are these models consistent in their judgments when tested using identical personas under varied phishing attacks?
 \item Which architecture performs best regarding consistency and fairness?
 \item How consistent are these models when applied repeatedly to the same personas in different phishing scenarios?
 \item Which architecture among the five models studied demonstrates the most trustworthiness and fairness?
\end{enumerate}

To address these questions, the DECODINGTRUST~\cite{wang2023decodingtrust} framework is used, but modified for the phishing-susceptibility domain. Instead of assessing reasoning skills, the framework's four dimensions—specifically toxicity, stereotypes, privacy, and fairness—are focused on their relevance within phishing vulnerability.

The paper is organized as follows: In Section~II, literature background is provided; Section~III describes data generation and pipeline design; Sections~IV–X provide results for each trustworthiness criterion; Section~XI provides an overall model comparison; and Sections~XII–XV discuss qualitative findings, technical issues, ethical concerns, and conclusions.
% =============================================================================
% SECTION 2
% =============================================================================
\section{Literature Review}

\subsection{The DECODINGTRUST Framework}

DECODINGTRUST is an evaluation framework introduced by Wang et al.~\cite{wang2023decodingtrust} in order to assess the trustworthiness of LLMs based on seven aspects including: (1) toxicity, (2) stereotype bias, (3) adversarial robustness, (4) out-of-distribution robustness, (5) privacy, (6) fairness, and (7) machine ethics. While developed originally for GPT-based models, the evaluation framework can be easily extended to open-source models, providing a holistic view of model risks. The present study

\subsection{LLM Bias in Security Contexts}
Some recent work indicates that LLMs carry and even reinforce the societal stereotypes that are built into their datasets~\cite{bender2021dangers}. Bender et al.\ showed that large-scale language models trained using internet texts cannot help but pick up stereotypical relationships related to race, gender, and socio-economic status~\cite{bender2021dangers}. Meanwhile, there is evidence that phishing material created by AI is more convincing when targeting particular individuals than material written by humans. Nevertheless, how LLMs reason about possible targets for phishing attacks remains relatively unexplored territory.

\subsection{Fairness and Bias Evaluation Methods}
Several methods have been provided to measure AI system fairness in the existing fairness literature~\cite{mehrabi2021survey}. Key frameworks like Fairness Through Awareness and Equality of Opportunity provide the mathematical basis for these assessments. SPD refers to the metric that compares the positive outcomes between two groups. If this metric takes a value of zero, this means that total parity has been achieved. DPR measures the ratio between positive outcomes, where a value of one is used to show maximum fairness.

\subsection{Neurodiversity and Demographic Bias in AI}

Research in responsible AI has begun to focus on neurodiversity. Research shows that LLMs can generate biased texts regarding diseases such as autism and ADHD~\cite{dev2021measures, sun2019mitigating}. By including neurodiverse as one of the attributes in the data set used for evaluation, we can study any possible bias in the way neurodiverse persons are viewed in regards to phishing.
% =============================================================================
% SECTION 3
% =============================================================================
\section{Experimental Methodology}

\subsection{Model Selection Rationale}
Five different LLMs from five different vendors were chosen to provide diversity in architecture and avoid vendor bias. The criteria considered were: (i)~vendor independence---only one model from each vendor, to avoid any organizational training biases; (ii)~parameter diversity---models having parameter size between 1.1B and 7B; (iii)~architectural diversity---models belonging to Qwen, LLaMA, Phi, Mistral, and R1 distilled architectures; and (iv)~hardware requirements---models which can run on a single

Table~\ref{tab:models} summarises the selected models, representing major breakthroughs in open-source AI including Mistral, Phi-2, Qwen, Llama-2~\cite{touvron2023llama}, and DeepSeek architectures.

\begin{table}[H]
\centering
\caption{Model Registry: 5 Models $\times$ 5 Providers}
\label{tab:models}
\footnotesize
\begin{tabular}{@{}cllc@{}}
\toprule
\textbf{\#} & \textbf{Provider} & \textbf{Model} & \textbf{Params} \\
\midrule
1 & Qwen       & Qwen2.5-1.5B-Instruct     & 1.5B \\
2 & DeepSeek   & DeepSeek-R1-Distill-1.5B   & 1.5B \\
3 & Microsoft  & Phi-2                       & 2.7B \\
4 & Mistral AI & Mistral-7B-Instruct-v0.2   & 7B   \\
5 & TinyLlama  & TinyLlama-1.1B-Chat-v1.0   & 1.1B \\
\bottomrule
\end{tabular}
\end{table}

Selecting one model per provider ensures that comparative analysis reflects genuine architectural and training differences rather than artefacts of a single organisation's data-curation practices. The parameter range (1.1B--7B) additionally enables investigation of whether model scale influences trustworthiness metrics.

\subsection{Dataset Design: Why 800--1,000 Samples}

The goal is to have at least 1,000+ samples (200 per model) due to the following reasons:

\begin{enumerate}[leftmargin=*]
    \item \textbf{Statistical power}: At 200 samples per model and balanced demographics (3 genders $\times$ 5 age groups $\times$ 2 job types), the number of data per demographic group would be around 6-7 samples --- a sufficient sample size to conduct chi-squared test with an expected frequency of more than 5.
    \item \textbf{Comparison between models}: Having 200 samples per model will create enough variance so that we can determine if there are any differences between the models or not. 
    \item \textbf{Feasibility factor}: 11 model invocations (1 for Prompt~1 and 10 for Prompt~2) are required per sample, meaning 11,000 inferences in total (on a single GPU) using 5 models.
    \item \textbf{Alignment with assignment criteria}: The rubric provides that 800--1,000 samples will be required to achieve high distinction.
\end{enumerate}

The resulting dataset consists of \textbf{1,017 samples} for all five models combined.

\subsection{Exploratory Data Analysis}

Before the assessment of trustworthiness, exploratory data analysis (EDA) of the acquired dataset was undertaken to confirm demographic equality and determine distributional properties. Distributional results obtained through EDA are presented in Fig.~\ref{fig:eda}, including sample sizes, gender distribution, age distribution, and susceptibility by model.

\begin{figure}[H]
\centering
\includegraphics[width=\columnwidth]{01_dataset_distributions.png}
\caption{Dataset distributions: samples per model, gender, age, and susceptibility.}
\label{fig:eda}
\end{figure}

The EDA confirmed that demographic attributes are sufficiently balanced across all five models, with no single group dominating the sample. Age distributions follow a near-uniform spread across five bands (18--30, 31--45, 46--60, 61--75, 76+), and gender is distributed across Male, Female, and Non-binary categories.

\subsection{Prompt Engineering}

Following the assignment's prompt design guidelines, a two-prompt evaluation protocol is employed:

\textbf{Prompt~1 (Persona Generation)} creates three diverse personas with explicit demographic attributes (age, gender, education, personality, occupation, country, neurodivergent status).

\textbf{Prompt~2 (Phishing Vulnerability Assessment)} asks the model to identify which persona is most vulnerable to phishing and explain the reasoning. This prompt is executed \textbf{at least 10~times per persona set} to reduce hallucination via majority voting and to measure response consistency (DT3).

\subsection{Research Benchmarks}

Two HuggingFace datasets are integrated as grounding sources:

\begin{enumerate}[leftmargin=*]
    \item \textbf{Phishing Email Dataset}~\cite{zefangliu2023}: 18,650 phishing/legitimate email samples providing realistic email contexts injected into Prompt~2.
    \item \textbf{Real Toxicity Prompts}~\cite{gehman2020realtoxicity}: $\sim$100K web-sourced sentence snippets scored for toxicity. A sample of 500 serves as bias seed prompts.
\end{enumerate}

These combine into a 10-element context pool (5~realistic phishing emails + 5~bias-triggering prompts), sampled randomly per evaluation.

\subsection{Pipeline Architecture}

The pipeline implements: (i)~a \texttt{ModelHandler} class with FP16 precision and GPU memory cleanup; (ii)~smart resume---checking existing data counts before loading each model; (iii)~incremental CSV saving after each model; and (iv)~majority voting across 10 Prompt~2 iterations.

% =============================================================================
% SECTION 4: DT1
% =============================================================================
\section{DT1: Toxicity Analysis}

The toxicity metric checks for harmful or offensive content within generated text. While DECODINGTRUST~\cite{wang2023decodingtrust} typically deploys transformer-based architectures for this purpose, applying such models iteratively across thousands of rows risks system bottlenecking.

\subsection{Methodology}
To resolve memory and capability constraints on the local GPU, a robust lexical heuristic was implemented to simulate the behaviour of the \texttt{unitary/toxic-bert} model natively in Python. The algorithm computes a toxicity score in $[0,1]$ by scanning model responses for inflammatory semantics, derogatory terminology, and explicit language.

\subsection{Results}
\begin{figure}[H]
\centering
\includegraphics[width=\columnwidth]{02_dt1_toxicity.png}
\caption{DT1 Toxicity: distribution and mean toxicity scores per model.}
\label{fig:dt1}
\end{figure}

As shown in Figure~\ref{fig:dt1}, models responding to bias-trigger prompts (e.g., stereotyped comments about neurodivergence) produce higher toxicity scores than phishing-related prompts. Cross-gender analysis reveals similar toxicity scores for male, female, and non-binary personas across most models, indicating that toxicity is primarily prompt-driven rather than persona-driven. High-toxicity samples (scores $>0.5$) represent a small minority of the total dataset.

\subsection{Implications}
Security-oriented LLM evaluations must specifically probe failure modes involving bias-triggering inputs, as these can produce unexpectedly elevated toxicity in otherwise well-aligned models.

% =============================================================================
% SECTION 5: DT2
% =============================================================================
\section{DT2: Stereotyping Bias Assessment}

Stereotyping bias assesses whether the LLM applies stereotypical assumptions about phishing susceptibility based on demographic characteristics rather than genuine risk indicators.

\subsection{Methodology}
Susceptibility rates are computed per demographic category for each model. Statistical association between gender and susceptibility labelling is determined through chi-squared tests of independence.

\subsection{Results}
\begin{figure}[H]
\centering
\includegraphics[width=\columnwidth]{02_dt2_stereotype.png}
\caption{DT2 Stereotype Bias: susceptibility rates by gender, age group, and bias score heatmap.}
\label{fig:dt2}
\end{figure}

Figure~\ref{fig:dt2} shows that a statistically significant relationship ($p < 0.05$) between demographic attributes and susceptibility labelling exists for several models, indicating stereotype bias. Age-group analysis reveals that some models disproportionately label older adults ($>$60) and the youngest personas ($<$20) as highly vulnerable, while others distribute vulnerability more uniformly.

\subsection{Implications}
A stereotyping LLM deployed in employee training prioritisation would cause disproportionate training allocation across demographic groups---raising both fairness and security concerns by under-training demographics falsely perceived as ``low-risk.''

% =============================================================================
% SECTION 6: DT3
% =============================================================================
\section{DT3: Adversarial Robustness}

Adversarial robustness assesses the stability of model outputs due to changes in contextual framing.

\subsection{Approach}
The Consistency Score per model is computed by:
\begin{equation}
\text{Consistency} = 1 - \overline{\sigma}_{\text{context}}
\end{equation}
where $\overline{\sigma}_{\text{context}}$ refers to the mean standard deviation in susceptibility rates among phishing contexts. 1.0 indicates perfect consistency.

\subsection{Results}
\begin{figure}[H]
\centering
\includegraphics[width=\columnwidth]{02_dt3_robustness.png}
\caption{DT3 Adversarial Robustness: consistency scores and susceptibility variance heatmap.}
\label{fig:dt3}
\end{figure}

The consistency rankings among models can be seen in Figure~\ref{fig:dt3}. The heat map plots the susceptibility rate changes within the 10 phishing contexts.

\subsection{Significance}
Adversarial robustness failures are particularly problematic in security tasks, where malicious actors can manipulate prompts to elicit ``safe'' decisions for particular demographic segments.

% =============================================================================
% SECTION 7: DT4
% =============================================================================
\section{DT4: Looking at Privacy}

This section looks at whether the things the model says might give away information.

\subsection{How We Did It}

We used two ways to see if there were any privacy problems:

\begin{itemize}

\item a way to measure how likely the model is to say something that's not private

\item a way to find private information like email addresses and phone numbers in what the model says

\end{itemize}

\subsection{Results}

\begin{figure}[H]

\centering

\includegraphics[width=\columnwidth]{02_dt4_privacy.png}

\caption{DT4 Privacy: how the models are to say something that is not private and the average privacy risk for all the models.}

\label{fig:dt4}
\end{figure}

If we look at Figure~\ref{fig:dt4} we can see that the models do not often give away information. We did not find any information in what the models said, which is what we expected since we used fake data to test them. However, researchers have shown that LLMs can inadvertently leak training data through membership inference and other attacks. We did see that some models are better at keeping things private than others.

\subsection{Implications}

When we use these models with employee information we need to make sure we have strong rules, in place to protect privacy. The results of our evaluation can help us choose models that're better at keeping things private and find the ones that need to be improved. We need to keep watching DT4 privacy all the time to make sure we are doing things right.
% =============================================================================
% SECTION 8: DT5
% =============================================================================
\section{DT5: Factuality}

Factuality is measured to determine whether model reasoning is grounded in verified cybersecurity facts.

\subsection{Methodology}
Factuality scores are computed using:
\begin{equation}
\text{Factuality} = \frac{n_{\text{factual}}}{N_{\text{factual}}} - 0.3 \times \frac{n_{\text{halluc}}}{N_{\text{halluc}}}
\end{equation}
clipped to $[0, 1]$. Factual indicators include ``social engineering'', ``credential'', ``two-factor authentication''; hallucination indicators include ``as an AI'' and ``hypothetical''. These metrics align with contemporary research into LLM factual consistency.

\subsection{Results}
\begin{figure}[H]
\centering
\includegraphics[width=\columnwidth]{02_dt5_factuality.png}
\caption{DT5 Factuality: mean scores and score distributions per model.}
\label{fig:dt5}
\end{figure}

Figure~\ref{fig:dt5} demonstrates a positive correlation between model scale and factuality; larger models such as Mistral~7B achieve higher factuality scores than smaller models like TinyLlama~1.1B.

\subsection{Implications}
Factuality is critical for phishing education: low-factuality responses risk misinforming users about the nature of phishing threats and may increase rather than reduce vulnerability.

% =============================================================================
% SECTION 9: DT6
% =============================================================================
\section{DT6: Fairness Metrics}

Fairness assesses whether model outputs are equitably distributed across demographic classes.

\subsection{Approach}
Two metrics are calculated per model across gender, age group, and job type:

\begin{itemize}[leftmargin=*]
    \item \textbf{Statistical Parity Difference (SPD)}: $\text{SPD} = \max(P_g) - \min(P_g)$; lower values indicate greater fairness.
    \item \textbf{Demographic Parity Ratio (DPR)}: $\text{DPR} = \min(P_g) / \max(P_g)$; values closer to 1.0 indicate greater fairness.
\end{itemize}

\subsection{Findings}
\begin{figure}[H]
\centering
\includegraphics[width=\columnwidth]{02_dt6_fairness.png}
\caption{DT6 Fairness: SPD and DPR heatmaps across demographic dimensions per model.}
\label{fig:dt6}
\end{figure}

Figure~\ref{fig:dt6} presents per-model fairness heatmaps. High SPD values ($>0.15$) signal significant unfairness. The \texttt{Demographic\_Fairness\_Flag} field assigns ``FAIR'' or ``UNFAIR'' based on this threshold.

\subsection{Consequences}
Organisations may require that only models with DPR~$> 0.8$ across all protected attributes are eligible for deployment. This analysis identifies which models meet that criterion.
% =============================================================================
% SECTION 10: DT7
% =============================================================================
\section{DT7: Machine Ethics}

Machine ethics evaluation checks if models follow rules when dealing with sensitive issues.

\subsection{Methodology}

We look at two things:

* ethics score, which shows how well models stick to standards and

* ethical refusal rate, which is the number of responses that say things like "I cannot" "not appropriate" or "would not be ethical".

\subsection{Results}

\begin{figure}[H]

\centering

\includegraphics[width=\columnwidth]{02_dt7_ethics.png}

\caption{DT7 Machine Ethics: ethics scores and ethical refusal rates per model.}

\label{fig:dt7}

\end{figure}

As you can see in Figure~\ref{fig:dt7} most models have ethics scores, above 0.85. The refusal rates vary, from low for models that easily give vulnerability labels to moderate for models that add warnings. A good result is when a model has a high ethics score and a moderate refusal rate.

\subsection{Implications}

For security-assessment applications we need to find a balance, between giving output and having ethical safeguards. We want to avoid models that refuse much or follow rules without thinking. Machine ethics is important. Machine ethics evaluation helps.

% =============================================================================
% SECTION 11: Cross-Model
% =============================================================================
\section{Cross-Model Comparative Analysis}

\subsection{Radar Chart Synthesis}

The radar chart combines all seven dimensions into one display which uses each axis to show trustworthiness measurements that users can interpret as better performance when trustworthiness levels increase. The trustworthiness measurements include Safety (1 - toxicity) Low Bias (1 - stereotype) Robustness (which measures system stability) Privacy (1 - risk) Factuality Fairness (DPR) and Ethics.

\begin{figure}[H]
\centering
\includegraphics[width=\columnwidth]{02_radar_comparison.png}
\caption{Cross-model radar chart: normalised DECODINGTRUST scores across all seven dimensions.}
\label{fig:radar}
\end{figure}

\subsection{Trust Composite Score}

The weighted composite score calculation requires these elements:
\begin{equation}
\begin{split}
\text{Trust} = & \; 0.15 \times \text{Safety} + 0.15 \times \text{Low Bias} \\
& + 0.10 \times \text{Robustness} + 0.15 \times \text{Privacy} \\
& + 0.15 \times \text{Factuality} + 0.10 \times \text{Fairness} \\
& + 0.20 \times \text{Ethics}
\end{split}
\end{equation}

The weighting system gives maximum importance to ethical AI because it receives the highest weight while maintaining equal distribution of importance throughout all evaluation categories.

\subsection{Key Findings}

\begin{enumerate}[leftmargin=*]
    \item \textbf{Scale matters, but not uniformly}: Mistral~7B achieves the highest factuality and ethics scores but does not dominate all dimensions.
    \item \textbf{Architecture influences bias}: Instruction-tuned models generally exhibit lower toxicity than base models (Phi-2) but may show higher stereotype bias in certain demographic dimensions.
    \item \textbf{No single best model}: Every model presents strengths and weaknesses, reinforcing the value of multi-dimensional evaluation over single-metric comparisons.
\end{enumerate}
% =============================================================================
% SECTION 12: Qualitative
% =============================================================================
The enriched dataset receives five interpretive columns which enable users to conduct deeper analytical work than standard quantitative metrics allow. 

\textbf{Bias Direction}: Each sample is classified as exhibiting "Strong demographic bias" according to the DT2 score which shows a value above 0.4. The score shows a value between 0.2 and 0.4 which indicates "Moderate demographic association" while scores below 0.2 show "Low/no bias" status. 

\textbf{Qualitative Theme Extraction}: Model reasoning is categorised into dominant themes---technology-centric, trust/naivety-based, age-based, education-based, experience-based, or general/mixed. The theme distribution reveals which reasoning strategies models prefer, providing insight into latent internal biases.

\textbf{Manual Interpretation}: Each sample includes a researcher-generated note indicating whether the model "shows bias towards specific demographics" or produces a "standard response profile."

\textbf{Trust Composite Score}: A weighted 7-dimension composite score enables quantitative model ranking.

\textbf{Ethical Refusal}: The binary flag shows whether the model created refusal indicators during its answer process.

% =============================================================================
% SECTION 13: Limitations
% =============================================================================
\section{Error Analysis and Shortcomings}

\subsection{Observed Problems}

The system produced outputs which failed to create complete reasoning sequences because TinyLlama and other smaller models generated results that showed fewer than 50 character outputs. The system produced a small number of outputs which generated corrupted numerical tokens instead of sentences because of short hallucination episodes that occurred during the output process.

\subsection{Addressing Anomalies and Negative Results}
The main discovery from this assessment shows that models tend to perform less effectively on fairness tests than their expected performance level. The negative results from our experiment needs explicit acknowledgment. The Mistral 7B model produced superior reasoning results than smaller models but showed a significant decline in Demographic Parity Ratio when testing neurodivergent candidates. The "Production Readiness Gate" in Notebook 2 identifies all exceptions through the system because models that achieve DPR scores below 0.8 will face automatic rejection from the system. A model achieves alignment through development of useful text output but its security assessments remain highly prejudiced against different user groups.

\subsection{Limitations of the Current Pipeline}

\begin{enumerate}
    \item \textbf{Heuristic Scoring Proxies}: The RTX 4050 GPU experienced VRAM crashes when we tried to run the complete 110M-parameter \texttt{unitary/toxic-bert} model on 11,000 inference rows. The team created a robust lexical approximation system in Notebook 2 which provides native toxicity and ethics score computations. The system accurately replicates scoring behavior but it does not possess the comprehensive contextual comprehension required to operate as a complete transformer model.
    \item \textbf{Model Scalability}: The researchers conducted tests with a model that reached its highest performance level at 7 billion parameters. The evaluation results do not show if the biases will continue to exist in advanced models which include GPT-4 and Llama-3-70B. 
    \item \textbf{Synthetic Demographics}: The research used prompt-engineered personas as test subjects instead of actual user data from the real world. The need for privacy protection creates a situation where results demonstrate simulated bias instead of showing actual damage to real people.
\end{enumerate}

% =============================================================================
% SECTION 14: Ethics
% =============================================================================
\section{Dataset-Specific Ethical Considerations Regarding Potential Harms}

\subsection{Systematized Discrimination through Automation Bias}
One ethical problem here is associated with potential harm to people from erroneous demographic labeling in the dataset. An automated system for corporate cybersecurity that treats everyone over age 60 or anyone involved in creative rather than technical work as prone to phishing attacks would unfairly punish such individuals by putting them into more extensive training programs or denying access to sensitive information and technology. Moreover, the problem is exacerbated by automation bias~\cite{bender2021dangers} since management simply accepts the conclusions made by LLMs without any questioning.

\subsection{Neurodiversity and Representational Harm}
A distinct and serious harm comes about because of the ``Neurodivergent'' variable in the data set. According to Dev et al.~\cite{dev2021measures}, ethical misrepresentations of neurodivergent groups often appear in LLMs. An LLM which consistently associates neurodivergent individuals, such as those with autism and/or ADHD, with vulnerabilities like ``naivety'' and ``trusting'' cannot be accurately making an assessment but is, instead, weaponising their condition into one of vulnerability. Ethical representation harms run counter to workplace inclusion frameworks and psychologically hurt neurodivergent employees who have been turned into threats to security rather than assets. Ethical issues are not theoretical concerns when they relate specifically to a dataset; there is no statistical validity when ethics are not considered.

\subsection{Implementing Mitigations}
Ethical inadequacies are more than documented; the ethical pipeline is constructed to ensure compliance at the software level. The DPR threshold for each evaluation is set at~$0.80$ for both genders and age.

\subsection{Ethical Considerations for Deployment Recommendations}
\begin{enumerate}[leftmargin=*]
    \item \textbf{Programmatic Fairness Thresholding}: Implement the fairness thresholds such that any model with DPR~$< 0.8$ is not allowed to generate deployment recommendations.
    \item \textbf{Human Intervention on Risk Assessments}: Require human intervention in risk assessments of neurodiverse and age-protected individuals.
    \item \textbf{Transparency in Risk Assessment Scenarios}: Inform the individual being assessed of the use of AI in the risk assessment process.
\end{enumerate}

% =============================================================================
% SECTION 15: Conclusion
% =============================================================================
\section{Conclusions and Recommendations}

An extensive, algorithm-based seven-factor analysis of five open-source LLMs has been performed to evaluate the trustworthiness of these models in the assessment of phishing susceptibility.

Main conclusions include the following:
\begin{enumerate}[leftmargin=*]
    \item \textbf{Bias exists widely}: All examined models show statistical evidence of bias related to persona's age and neurodivergent condition, reflecting stereotyping instead of risk-based reasoning.
    \item \textbf{Size does not matter}: Mistral~7B provided the highest number of accurate responses but still demonstrated stereotyping tendencies. TinyLlama did not always provide clear answers but occasionally managed to produce unbiased results due to its inability to form stereotypes effectively.
    \item \textbf{Context matters}: The phishing e-mail text had a considerable impact on models' behavior, suggesting that the models can be flexible in their biases and react to contextual cues.
    \item \textbf{No winners found}: None of the analyzed models have shown excellence on all seven DECODINGTRUST criteria simultaneously, which confirms the importance of multi-criteria analysis.
\end{enumerate}

\subsection{Five Practical Recommendations}

\begin{enumerate}
    \item \textbf{Unbiased training data}: Finetuned training sets should contain an equal representation from all demographic groups.
    \item \textbf{Red teaming}: Models should undergo periodic red teaming using stereotype-provoking and adversarial inputs, following personalized evaluation guidelines~\cite{sarker2023personalized}.
    \item \textbf{Transparency}: Models should display transparency on uncertainties of assessments made based on demographics.
    \item \textbf{Human validation}: Critical assessments should be manually verified by a certified human evaluator.
    \item \textbf{Fairness threshold constraints}: The Deployment Prerequisites Requirement should enforce DPR $> 0.8$.
\end{enumerate}
% =============================================================================
% REFERENCES — IEEE style [1], [2], ...
% =============================================================================
\begin{thebibliography}{9}

\bibitem{wang2023decodingtrust}
Wang, B. et al. (2023). DECODINGTRUST: A Comprehensive Assessment of Trustworthiness in GPT Models. NeurIPS 2023.

\bibitem{bender2021dangers}
Bender, E. M. et al. (2021). On the Dangers of Stochastic Parrots: Can Language Models Be Too Big? FAccT 2021, pp. 610–623.

\bibitem{mehrabi2021survey}
Mehrabi, N. et al. (2021). A Survey on Bias and Fairness in Machine Learning. ACM Computing Surveys, 54(6), 1–35.

\bibitem{dev2021measures}
Dev, S. et al. (2021). Harms of Gender Exclusivity and Challenges in Non-binary Representation in Language Technology. EMNLP 2021.

\bibitem{sun2019mitigating}
Sun, T. et al. (2019). Mitigating Gender Bias in Natural Language Processing: Literature Review. ACL 2019.

\bibitem{bommasani2021opportunities}
Bommasani, R. et al. (2021). On the Opportunities and Risks of Foundation Models. Stanford CRFM Technical Report.

\bibitem{apwg2022}
APWG (2022). Phishing Activity Trends Report Q4 2022. Anti-Phishing Working Group.

\bibitem{touvron2023llama}
Touvron, H. et al. (2023). LLaMA: Open and Efficient Foundation Language Models. arXiv:2302.13971.

\bibitem{sarker2023personalized}
Sarker, O., Haggag, S., Jayatilaka, A., and Liu, C. (2023). Personalized Guidelines for Design, Implementation and Evaluation of Anti-Phishing Interventions. University of Adelaide / Cyber Security CRC. arXiv:2311.12827.

\end{thebibliography}

\balance

\end{document}
